\section{What does a pattern promise?}\label{ch01:sec:pattern}
Suppose two people each have an odd number of tokens. If they put their tokens together, must the total be even? Try a few cases: $3+5=8$, $7+9=16$, and $1+1=2$. Every total is even. But there are infinitely many pairs of odd numbers, and no list of examples can include them all. The examples tell us what happened for the pairs we tried. They do not explain what must happen for a pair we have not tried.

To turn the pattern into mathematics, we first say exactly which numbers it is about. An \emph{integer} is a whole number that may be positive, negative, or zero: $\ldots,-2,-1,0,1,2,\ldots$. The collection of all integers is written $\ZZ$. Now the pattern becomes a precise statement:
\begin{center}
\emph{For every pair of odd integers, their sum is even.}
\end{center}
This statement says more than the token story. Nobody holds $-3$ tokens, but $-3$ is an odd integer, so the statement also covers sums such as $-3+5=2$ and $-7+(-1)=-8$. A calculation checks one particular pair. A proof explains why every pair the statement allows must work.

So when you bring a claim like this to a classmate or to an AI assistant, the first question to settle is not ``Can you prove this?'' It is ``Which numbers does the claim talk about, and what exactly must be true of them?'' Suppose an assistant quietly replaces ``integers'' by ``positive integers'' and then gives a correct argument. It has answered a different question. The new statement happens to be true, but its proof says nothing about a sum such as $-3+5$.
